\chapter*{How to Use This Book}
\addcontentsline{toc}{chapter}{How to Use This Book}
The chapters are written so that the reader can reproduce the examples directly. The following elements are used consistently:
\begin{description}[style=nextline,leftmargin=32mm]
	\item[Introduction] Introduces the examples
	\item[Remark] Adds interpretation, limitations or practical advice.
	\item[Example] Presents a worked application.
	\item[Calculation] Aligns formulas, substituted values and results.
	\item[Simulation] Describes model settings and simulation steps.
	\item[Chapter Summary] Collects the main conclusions at the end of a chapter.
\end{description}
 
\vspace{11mm}

It is recommended that you work through the book from beginning to end, while reproducing the examples in LTspice. A prerequisite for successful simulations is a good understanding of how to use the LTspice simulator. A solid understanding of SPICE syntax is absolutely essential for a thorough understanding of how the models are constructed.
The complex models of the transformer and the synchronous generator are explained in the appendices of this book. To better understand the detailed modeling, it is recommended to review the LTspice-library files.

The free software WxMaxima was used in the book for simple calculations. The free software GeoGebra was used to construct the vector diagrams. In the “coupled lines” section, the Atp (Atpdraw) program was used for the overhead line model; however, it is not required for the simulation.

%\cleardoublepage